% Originally: the \section{Solutions} of content/week-03.tex

\section{Solutions}
\label{sec:connectedness_solutions}

% Transition: Claude Opus 5 (High)
Not every solution in this chapter is collected here. \Cref{ex:connectedness_tf}
(\cpageref{ex:connectedness_tf}) and \cref{ex:disk_connected} (\cpageref{ex:disk_connected}) are
solved where they appear.

\begin{exercisesolution}[Proof of \cref{ex:ai_connected_image}]
% Generator: Claude Opus 5 (High)
\begin{enumerate}[label=\textbf{(\alph*)}]
  \item Let $p, q \in f(X)$, say $p = f(x)$ and $q = f(y)$. Since $X$ is path-connected there is
    a path $\gamma : [0,1] \to X$ with $\gamma(0)=x$ and $\gamma(1)=y$. Then
    $f \circ \gamma : [0,1] \to f(X)$ is continuous as a composition of continuous maps, and it
    satisfies $(f\circ\gamma)(0) = p$, $(f\circ\gamma)(1) = q$. So it is a path in $f(X)$ from
    $p$ to $q$, and $f(X)$ is path-connected.
  \item A convex set $U$ is path-connected: the straight segment $t \mapsto (1-t)x + ty$ lies in
    $U$ and is continuous. The map $\varphi : U \to \mathbb{R}^n\times\mathbb{R}^m$,
    $\varphi(x) := (x,g(x))$, is continuous because each component is, and
    $\Gamma_g = \varphi(U)$. Apply \textbf{(a)}.
\end{enumerate}

\begin{remark}
Part \textbf{(a)} is strictly easier than its connectedness counterpart
(\cref{thm:continuity_preserves_connected}), because path-connectedness is a statement one
verifies by \emph{exhibiting} something, namely a path, and continuous maps push paths forward
directly. Connectedness has no such witness to push forward, which is why its proof has to run
by contradiction through preimages instead. This asymmetry is worth noticing: it is the reason
path-connectedness is usually the easier of the two to prove, and connectedness the easier of
the two to \emph{use}.
\end{remark}
\end{exercisesolution}

% ===========================================================================
% Official sheet 3 (3.1--3.11). Derived independently, then checked against
% exercises/Sol3_Analysis2_eng.pdf.
% ===========================================================================

\begin{exercisesolution}[Proof of \cref{ex:4.2}]
% Generator: Claude Sonnet 5 (Medium)
Consider the map
\[
\varphi : U \to \mathbb{R}^n \times \mathbb{R}^m, \qquad x \mapsto (x, f(x)).
\]
Since $f \in C^1(U,\mathbb{R}^m)$ is in particular continuous, and the identity on $U$ is
continuous, $\varphi$ is continuous. By definition $\Gamma_f = \varphi(U)$. Since $U$ is connected
and $\varphi$ is continuous, \cref{thm:continuity_preserves_connected} gives that $\varphi(U) =
\Gamma_f$ is connected.
\end{exercisesolution}

\begin{exercisesolution}[Solution of \cref{ex:closed_unbounded_not_simply_connected}]
% Generator: Claude Opus 5 (High)
Take the closed exterior of the unit disc,
\[U := \{x \in \mathbb{R}^2 \mid \lvert x\rvert \ge 1\}.\]
\begin{itemize}
  \item \textbf{Closed:} $U$ is the preimage of the closed set $[1,\infty)$ under the continuous map
    $x \mapsto \lvert x\rvert$.
  \item \textbf{Unbounded:} it contains $(k,0)$ for every $k \in \mathbb{N}$.
  \item \textbf{Connected:} it is in fact path-connected. Given $p,q \in U$, travel radially outwards
    from $p$ to the circle of radius $R := \max\{\lvert p\rvert, \lvert q\rvert\}$, then along that
    circle, then radially inwards to $q$. Every point on this path has norm between $1$ and $R$, so
    the path stays in $U$.
  \item \textbf{Not simply connected:} the loop $\gamma(t) := (2\cos 2\pi t,\, 2\sin 2\pi t)$ winds
    once around the origin, and no loop with winding number $1$ about a point outside $U$ can be
    contracted to a point within $U$. Contracting it would force the loop to sweep across the
    removed open disc $\{\lvert x\rvert < 1\}$.
\end{itemize}
\begin{ainote}
The exercise is really a test of whether the three notions have been kept apart. Closed-and-unbounded
rules out the reflex answer of an annulus $\{1 \le \lvert x\rvert \le 2\}$, and \qt{connected but not
simply connected} rules out a half-plane or all of $\mathbb{R}^2$. Deleting an \emph{open} disc
rather than a point is what keeps the result closed while still puncturing it.
\end{ainote}
\end{exercisesolution}

\begin{exercisesolution}[Proof of \cref{ex:sphere_path_connected}]
% Generator: Claude Opus 5 (High)
Both directions are proved by looking at what $S^{n-1}$ actually is in the two cases.

\textbf{The case $n = 1$.} Here $S^0 = \{x \in \mathbb{R} \mid x^2 = 1\} = \{-1, +1\}$, a two-point
space. Suppose $\gamma : [0,1] \to S^0$ were a path with $\gamma(0) = -1$ and $\gamma(1) = 1$.
Regarded as a map into $\mathbb{R}$, it is continuous on an interval, so by the intermediate value
theorem it takes the value $0$ somewhere. But $0 \notin S^0$, a contradiction. Hence $S^0$ is not
path-connected.

\textbf{The case $n \ge 2$.} Fix $p, q \in S^{n-1}$, and suppose first that $q \neq -p$. Then the
straight segment between them misses the origin: if $(1-t)p + tq = 0$ for some $t \in [0,1]$, taking
norms gives $(1-t) = t$, so $t = \tfrac12$ and $p = -q$, which we excluded. We may therefore
normalise the segment and set
\[\gamma(t) := \frac{(1-t)p + tq}{\lvert (1-t)p + tq\rvert}, \qquad t \in [0,1].\]
This is continuous, takes values in $S^{n-1}$, and satisfies $\gamma(0) = p$, $\gamma(1) = q$.

It remains to handle $q = -p$, and this is exactly where $n \ge 2$ enters: the sphere $S^{n-1}$ then
contains a point $r$ with $r \neq \pm p$. (For $n \ge 2$ the sphere is infinite --- pick any unit
vector $v$ orthogonal to $p$, which exists because $p^\perp$ has dimension $n - 1 \ge 1$.) Neither
$r$ nor $q$ is the antipode of the other's partner, so the previous construction gives a path from
$p$ to $r$ and a path from $r$ to $q$; concatenating them joins $p$ to $q$. Hence $S^{n-1}$ is
path-connected.
\begin{remark}
The construction in the main case is the radial projection of the segment, and it is the standard way
to move around a sphere: it works verbatim on any $S^{n-1}$ and explains why the antipodal pair is
the only obstruction. Note also that $n \ge 2$ is used only once, to produce a third point --- which
is precisely what $S^0$ does not have.
\end{remark}
\begin{ainote}
The official solution introduces the same construction with the words \qt{now let $n \ge 1$}, and
then disposes of the antipodal case $x = -y$ by composing two paths through \qt{an arbitrary other
point $w \in S^{n-1}$}. For $n = 1$ there is no such $w$: $S^0 = \{\pm1\}$ consists of $x$ and $-x$
and nothing else. The bound is therefore $n \ge 2$, exactly as the statement being proved requires
--- and as the official solution's own preceding paragraph, which shows $S^0$ is \emph{not}
path-connected, confirms. The step is written out explicitly above rather than left to the reader,
since it is the only place the hypothesis is used.
\end{ainote}
\end{exercisesolution}

\begin{exercisesolution}[Proof of \cref{ex:examHS24_TF_topology}]
% Generator: Gemini 3.1 Pro (High)
\begin{enumerate}[label=\textbf{(\alph*)}]
  \item \textbf{True:} $U$ is defined by non-strict inequalities involving continuous functions (specifically, the pre-image of $[0, \infty)$ under $g_1(x,y) = y+x^2$ and $g_2(x,y) = x^2-y$). Consequently, $U$ is closed.
  \item \textbf{True:} We can contract $U$ to the origin via the homotopy $H((x,y), t) = (tx, t^2y)$ for $t \in [0,1]$. Note that $(tx)^2 = t^2x^2$, and since $-x^2 \le y \le x^2$, multiplying by $t^2 \ge 0$ gives $-(tx)^2 \le t^2y \le (tx)^2$, which ensures the contraction stays within $U$.
  \item \textbf{False:} $U$ is not bounded (for example, the sequence $(n,0)$ is in $U$ for all $n \in \mathbb{N}$), and therefore not compact.
  \item \textbf{True:} Since $U$ is path-connected (from the homotopy to the origin) and therefore connected, its continuous image $f(U)$ is connected in $\mathbb{R}$.
  \item \textbf{False:} The function is not bounded on $U$. Consider the path $(x,0)$ inside $U$; here $f(x,0) = x^2$, which grows without bound as $x \to \infty$.
  \item \textbf{True:} Since $f(x,y) \ge 0$ everywhere on $U$ and $f(0,0) = 0$ with $(0,0) \in U$, the infimum is $0$ and it is attained at the origin.
\end{enumerate}
\end{exercisesolution}

\begin{exercisesolution}[Solution to \cref{ex:box_intersection_disconnected_union_connected}]
% Generator: Claude Opus 5 (High) --- the construction below is the official key's, with the
% strip widened from half-width 1/10 to 1/4 so that the figure is legible; every verification is
% written for these notes.
Take an open annulus and an open vertical strip through its hole:
\[U_1 := B_1(0) \setminus \overline{B_{1/2}(0)}
  = \big\{p \in \mathbb{R}^2 \ \big|\ \tfrac12 < \lvert p\rvert < 1\big\}, \qquad
  U_2 := \big(-\tfrac14,\tfrac14\big)\times\mathbb{R}.\]

% Generator: Claude Opus 5 (High) --- FIG-MOCK-19
% Geometry: two concentric circles of radii 1 and 1/2 (drawn at 2.2 cm and 1.1 cm per unit),
% and a vertical strip |x| < 1/4 (drawn at +-0.55 cm) running past both of them. All four
% boundaries are DASHED, because both sets are open and neither contains its boundary. The
% annulus is shaded blue, the strip orange, and the two pieces of the intersection purple on
% top of both. The purple pieces are two lenses, one above and one below the x-axis, and the
% whole point of the picture is that they do not touch.
% Arithmetic: on U_1 \cap U_2 we have |x| < 1/4 and |p| > 1/2, so
%   y^2 > 1/4 - x^2 >= 1/4 - 1/16 = 3/16,  i.e. |y| > sqrt(3)/4 = 0.433.
% At the drawing scale that is 0.433 * 2.2 = 0.953 cm, which is where each lens must stop:
% it is sqrt(1.1^2 - 0.55^2) = 1.1 * sqrt(3)/2 = 0.953 cm, the height at which the strip edge
% x = 0.55 crosses the inner circle. The two agree, as they must.
\begin{center}
\begin{tikzpicture}[>=Latex]
  % --- U_1, the annulus: fill the outer disc, then punch out the inner one ---
  \fill[blue!12] (0,0) circle (2.2);
  \fill[white] (0,0) circle (1.1);
  % --- U_2, the strip (covers the annulus shading where they meet) ---
  \fill[orange!14] (-0.55,-2.9) rectangle (0.55,2.9);
  % --- the intersection, painted back on top of both: purple over the outer disc inside the
  %     strip, then orange restored over the hole, which belongs to U_2 alone ---
  \begin{scope}
    \clip (-0.55,-2.9) rectangle (0.55,2.9);
    \fill[purple!40] (0,0) circle (2.2);
    \fill[orange!14] (0,0) circle (1.1);
  \end{scope}
  % --- boundaries, all dashed: both sets are open ---
  \draw[blue!70!black, dashed, thick] (0,0) circle (2.2);
  \draw[blue!70!black, dashed, thick] (0,0) circle (1.1);
  \draw[orange!85!black, dashed, thick] (-0.55,-2.9) -- (-0.55,2.9);
  \draw[orange!85!black, dashed, thick] (0.55,-2.9) -- (0.55,2.9);
  % --- labels ---
  % label placed inside the blue band: radius 1.65 at 135 degrees, between 1.1 and 2.2
  \node[blue!70!black, font=\small] at (-1.17,1.17) {$U_1$};
  \node[orange!85!black, font=\small, anchor=south] at (0,2.95) {$U_2$};
  \node[purple!75!black, font=\small, align=left, anchor=west] at (2.7,0)
    {$U_1\cap U_2$:\\two components};
  \draw[purple!75!black, ->] (2.65,0.35) to[out=150,in=20] (0.62,1.6);
  \draw[purple!75!black, ->] (2.65,-0.35) to[out=210,in=340] (0.62,-1.6);
  \draw[gray, dotted] (-2.6,0) -- (2.4,0);
\end{tikzpicture}
\end{center}

Four things have to be checked, and the figure does none of them by itself.

\begin{enumerate}[label=\textbf{(\alph*)}]
  \item \textbf{$U_1$ is open and connected.} It is the intersection of the open set $B_1(0)$
    with the complement of the closed set $\overline{B_{1/2}(0)}$, hence open. It is
    path-connected, since any two of its points can be joined by moving radially to the circle of
    radius $\tfrac34$ and then along that circle, and a path-connected set is connected.
  \item \textbf{$U_2$ is open and connected.} It is a product of an open interval with
    $\mathbb{R}$, hence open, and it is convex, hence path-connected and so connected.
  \item \textbf{$U_1\cap U_2$ is not connected.} If $(x,y)$ lies in it then $\lvert x\rvert <
    \tfrac14$ and $x^2+y^2 > \tfrac14$, so
    \[y^2 > \tfrac14 - x^2 > \tfrac14 - \tfrac1{16} = \tfrac3{16} > 0 .\]
    In particular $y \neq 0$ throughout, so the two open half-planes $\{y>0\}$ and $\{y<0\}$
    separate $U_1\cap U_2$. Both pieces are non-empty, since $(0,\tfrac34)$ and $(0,-\tfrac34)$
    both lie in $U_1\cap U_2$. That is a separation, so the intersection is disconnected.
  \item \textbf{$U_1\cup U_2$ is connected.} Both sets are connected and they share the point
    $(0,\tfrac34)$, so \cref{cor:union_connected_common_point} applies to the two-element family
    $\{U_1,U_2\}$.
\end{enumerate}

\begin{remark}[Why the intersection is the odd one out]
\Cref{cor:union_connected_common_point} makes unions easy, and nothing in this chapter says
anything at all about intersections. That silence is not an oversight. A union of connected sets
sharing a point cannot be split, because any separation would have to split one of the pieces;
an intersection carries no such argument, and the picture above shows why: the strip enters the
annulus twice, and what it takes away in between is exactly what would have joined the two
pieces.

Note also that the construction needs the hole. If $U_1$ were a disc, then $U_1 \cap U_2$ would
be convex and the example would collapse; the failure of connectedness in \textbf{(c)} is the
annulus's non-simple-connectedness showing up in a second guise.
\end{remark}
\end{exercisesolution}
